\documentclass[11pt,a4paper]{article}

\usepackage[T1]{fontenc}
\usepackage[utf8]{inputenc}
\IfFileExists{lmodern.sty}{\usepackage{lmodern}}{}
\usepackage[margin=2.5cm]{geometry}
\usepackage{amsmath,amssymb,amsthm}
\usepackage{booktabs}
\usepackage{graphicx}
\usepackage{xcolor}
\usepackage{listings}
\usepackage[colorlinks=true,linkcolor=blue!50!black,urlcolor=blue!50!black,
            citecolor=blue!50!black]{hyperref}

\lstset{basicstyle=\ttfamily\footnotesize, breaklines=true, frame=single,
        framesep=4pt, rulecolor=\color{black!30},
        backgroundcolor=\color{black!3}, columns=fullflexible}

\theoremstyle{definition}
\newtheorem{definition}{Definition}[section]
\newtheorem{assumption}{Assumption}[section]
\theoremstyle{plain}
\newtheorem{proposition}{Proposition}[section]
\theoremstyle{remark}
\newtheorem{remark}{Remark}[section]

\newcommand{\Neff}{N_{\mathrm{eff}}}
\newcommand{\Ind}{\mathbf{1}}
\newcommand{\Rr}{\mathcal{R}}
\newcommand{\Tt}{\mathcal{T}}
\newcommand{\Prob}{\mathbb{P}}
\newcommand{\Exp}{\mathbb{E}}
\newcommand{\Var}{\mathrm{Var}}
\newcommand{\dif}{\mathrm{d}}

\title{\bfseries A scale-invariant, distribution-free procedure for comparing\\
weighted muon fluxes from different simulation productions}
\author{Muon flux validation}
\date{\today}

\begin{document}
\maketitle

\begin{abstract}
\noindent
We define and validate a procedure for testing whether the muon fluxes
delivered by two Monte Carlo productions are statistically compatible. The
procedure is constrained by three features of the data: the total weight of a
sample is a generation setting and carries no physical information, so all
statistics must be invariant under independent rescaling of either sample's
weights; the distributions are non-Gaussian with heavy tails that carry
physics; and the events are weighted with strongly concentrated weights, so
the number of muons is not the number of independent observations (Sec.~\ref{sec:neff}). We adopt
(i) ratios of weighted sums within a single sample, with uncertainties from a
nonparametric bootstrap over resampling units, and (ii) a weighted two-sample
Anderson--Darling statistic whose null distribution is obtained by permutation
of unit labels. We state the assumptions under which each is valid, give the
exact discrete forms evaluated by the implementation, and quantify the
approximations introduced by binning and by the substitution of effective
sample sizes for counts. We show that permutation calibration renders the
validity of the test independent of both approximations, which affect power
only. The procedure is validated by unit tests, closure, a calibration scan of
the p-value distribution under the null, a binning-stability scan, and an audit
of the weight bookkeeping of split productions (Sec.~\ref{sec:audit}). We report
its application to a pair of FairShip productions.
\end{abstract}

\tableofcontents

% ===========================================================================
\section{Introduction and problem statement}

Let two Monte Carlo productions each deliver a set of muons crossing a fixed
reference plane, each muon carrying a statistical weight. We wish to test the
hypothesis that the two productions generate muons from the same distribution
over the observable phase space.

Section~\ref{sec:setup} fixes the notation and states the null hypothesis;
Sec.~\ref{sec:scale} states the invariance the data impose. Four properties of
the data determine the method.

\begin{enumerate}
\item[(P1)] \textbf{The normalisation is not observable.} The productions are
generated with deliberately different muon rates; the total weight of a sample
is therefore a configuration choice. Any statistic sensitive to it measures the
generator setup rather than the physics. Since the normalisation difference is
typically the largest single discrepancy in the data, this is a binding
constraint rather than a technicality.

\item[(P2)] \textbf{The distributions are not Gaussian.} Momentum spectra fall
over several orders of magnitude and angular distributions are sharply peaked.
Procedures that compare only first and second moments --- for instance a Fisher
linear discriminant --- have zero power against alternatives that preserve
those moments. Section~\ref{sec:validation} exhibits such an alternative
explicitly.

\item[(P3)] \textbf{The tails carry information.} The physics of interest
includes the high-momentum and large-angle regions, which are sparsely
populated. A test must weight discrepancies there commensurately with their
statistical precision rather than being dominated by the bulk.

\item[(P4)] \textbf{The observations are weighted and correlated.} Biased
generation produces highly non-uniform weights; muons originating in the same
primary interaction, and copies of the same muon produced by variance-reduction
splitting, are mutually dependent.
\end{enumerate}

% ===========================================================================
\section{Notation and formal setup}
\label{sec:setup}

\begin{definition}[Sample]
\label{def:sample}
A \emph{sample} is a finite multiset
\begin{equation}
  \mathcal{S} = \bigl\{(x_i,\, w_i,\, u_i)\bigr\}_{i=1}^{n},
\end{equation}
where $x_i \in \mathbb{R}^d$ is the observable vector of muon $i$, $w_i > 0$ is
its weight, and $u_i \in \{1,\dots,m\}$ labels the \emph{resampling unit}
containing it. We write $\Rr$ for the reference sample (of size $n_\Rr$ over
$m_\Rr$ units) and $\Tt$ for the test sample.
\end{definition}

In Definition~\ref{def:sample} the weight $w_i$ is the value stored by the simulation for the track producing
the muon (in \textsc{FairShip}, \texttt{MCTrack.fW}), inherited by a plane hit
from the track that generated it. Its absolute normalisation is convention; see
(P1).

\begin{definition}[Resampling unit]
\label{def:unit}
A resampling unit is a subset of muons treated as a single indivisible
observation. Two choices are used: a \emph{file}, i.e.\ all muons written by
one simulation job; and an \emph{event block}, a set of whole tree entries
assigned by a fixed hash. The unit is the file when $m \ge 5$ and an event
block otherwise. Every muon belongs to exactly one unit, and units partition
the sample.
\end{definition}

The purpose of Definition~\ref{def:unit} is to make the objects resampled by the
permutation and bootstrap procedures mutually independent, or close enough to
it, despite (P4): muons sharing a primary interaction, and splitting copies of a
single muon, lie in the same tree entry and therefore in the same unit.

\begin{definition}[Observables]
For a muon with momentum $(p_x,p_y,p_z)$, transverse position $(x,y)$ at the
plane, and electric charge $q \in \{+1,-1\}$, write
$p = (p_x^2+p_y^2+p_z^2)^{1/2}$, $p_T = (p_x^2+p_y^2)^{1/2}$ and
$\theta = \arctan\!\left(p_T/|p_z|\right)$. The eight scalar observables tested
are
\begin{align}
  x^{(1)} &= \log_{10}(p/\mathrm{GeV}), &
  x^{(2)} &= \log_{10}(p_z/\mathrm{GeV}), \notag\\
  x^{(3)} &= \operatorname{arcsinh}\!\left(p_T/\varepsilon\right),\
            \varepsilon = 0.05~\mathrm{GeV}, &
  x^{(4)} &= \log_{10}(\theta/\mathrm{mrad}), \notag\\
  x^{(5)} &= x, \quad x^{(6)} = y, &
  x^{(7)} &= q\,x, \quad x^{(8)} = (x^2+y^2)^{1/2}. \label{eq:obs}
\end{align}
Only quantities measured at the plane enter; no generator-level or ancestry
information is used. Muons are selected by $|\mathrm{PDG}| = 13$,
$p_{\min} \le p \le p_{\max}$ and $p_z > 0$.
\end{definition}

The transformations in Eq.~\eqref{eq:obs} are monotone reparametrisations. The
Anderson--Darling statistic defined in Sec.~\ref{sec:ad} is invariant under
strictly monotone transformations of the observable, so the choice of
$\log_{10}$ or $\operatorname{arcsinh}$ affects only the binning and the
readability of the plots, not the test itself. The variable
$x^{(7)} = q\,x$ is included because in a magnetised absorber the sign of the
horizontal displacement relative to the charge carries information that $|x|$
discards.

\begin{definition}[Null hypothesis]
\label{def:h0}
Let $\Rr$ and $\Tt$ be labelled collections of units. The null hypothesis
$H_0$ is that the units are \emph{exchangeable} between the two labels: the
joint law of the observed data is invariant under any permutation of the
assignment of units to $\{\Rr, \Tt\}$ that preserves the group sizes
$(m_\Rr, m_\Tt)$.
\end{definition}

Definition~\ref{def:h0} is stronger than ``the two marginal distributions
coincide'' and weaker than ``all muons are i.i.d.'': it permits arbitrary
dependence within a unit. It is the hypothesis that the permutation test of
Sec.~\ref{sec:perm} tests exactly.

% ===========================================================================
\section{Scale invariance}
\label{sec:scale}

\begin{assumption}[Arbitrary normalisation]
\label{ass:scale}
The physical content of a sample is unchanged under $w_i \mapsto a\,w_i$ for any
$a>0$, independently for $\Rr$ and $\Tt$.
\end{assumption}

Assumption~\ref{ass:scale} is a statement about the generation procedure, not
about the data; it holds because the muon rate per proton on target is set by
the production configuration. It requires every reported statistic $S$ to
satisfy
\begin{equation}
  S\bigl(\{a w^{\Rr}\}, \{b w^{\Tt}\}\bigr)
  = S\bigl(\{w^{\Rr}\}, \{w^{\Tt}\}\bigr)
  \qquad \forall\, a,b>0.
  \label{eq:invariance}
\end{equation}

Two consequences follow. First, every reported quantity is a ratio of weighted
sums \emph{within one sample} (Sec.~\ref{sec:fractions}), which satisfies
Eq.~\eqref{eq:invariance} identically. Second, the shape statistic is computed
from normalised empirical distribution functions, and each permutation replica
is renormalised in the same way, so that no relabelling can respond to the
normalisation difference (Sec.~\ref{sec:perm}).

\begin{remark}
Equation~\eqref{eq:invariance} is a genuine restriction and not automatically
satisfied by standard practice. A classifier trained to discriminate $\Rr$ from
$\Tt$ using raw weights will learn the normalisation offset, which is the
largest and most easily learned feature present, and report a separation with no
physical content. The implementation enforces the invariance by test: one
sample's weights are multiplied by $10^{3}$ and bit-identical statistics and
p-values are required.
\end{remark}

% ===========================================================================
\section{Effective sample size}
\label{sec:neff}

\begin{definition}[Effective sample size]
For a sample with weights $\{w_i\}_{i=1}^n$,
\begin{equation}
  \Neff = \frac{\left(\sum_{i=1}^{n} w_i\right)^{2}}{\sum_{i=1}^{n} w_i^{2}} .
  \label{eq:neff}
\end{equation}
\end{definition}

Equation~\eqref{eq:neff} is Kish's effective sample size~\cite{Kish1965}. Its
justification is as follows. Let $\bar{y} = \sum_i w_i y_i / \sum_i w_i$ be a
weighted mean of independent quantities $y_i$ with common variance $\sigma^2$,
the weights treated as fixed. Then
\begin{equation}
  \Var(\bar{y})
  = \sigma^{2}\,\frac{\sum_i w_i^{2}}{\left(\sum_i w_i\right)^{2}}
  = \frac{\sigma^{2}}{\Neff},
\end{equation}
so $\Neff$ is the number of equally weighted observations giving the same
variance. The ratio $n/\Neff \ge 1$ is the design effect, with equality if and
only if all weights are equal.

\begin{proposition}[Inflation of $\Neff$ by splitting]
\label{prop:split}
Suppose each muon of weight $w$ is replaced by $k$ copies of weight $w/k$.
Then $\sum_i w_i$ is unchanged, $\sum_i w_i^{2}$ is divided by $k$, and hence
$\Neff$ is multiplied by $k$.
\end{proposition}

\begin{proof}
$k \cdot (w/k) = w$ and $k \cdot (w/k)^2 = w^2/k$; substitute in
Eq.~\eqref{eq:neff}.
\end{proof}

Proposition~\ref{prop:split} implies that for a production employing splitting,
Eq.~\eqref{eq:neff} is an \emph{upper bound} on the effective number of
independent observations, attained only if the copies are statistically
independent. Since copies of one muon are propagated separately through the
absorber they are neither identical nor independent, and the true value lies
between $\Neff/k$ and $\Neff$. Reported values of $\Neff$ for split samples must
be read with this in mind.

The uncertainties quoted in this work do not use Eq.~\eqref{eq:neff}: they are
obtained by resampling whole units (Sec.~\ref{sec:bootstrap}), under which
splitting copies move together and their dependence is accounted for
automatically. Empirically, for the samples of Sec.~\ref{sec:results} the ratio
of bootstrap variances between a split and an unsplit sample was $1.3$, against
a factor $15$ implied by the ratio of their $\Neff$ values --- confirming both
that Eq.~\eqref{eq:neff} overstates the information content and that the
bootstrap does not inherit the error.

% ===========================================================================
\section{Tests of relative amounts}
\label{sec:fractions}

\subsection{Estimands and estimators}

\begin{definition}[Flux fraction]
For a measurable region $R \subseteq \mathbb{R}^d$ of muon phase space, the
flux fraction of a sample is the estimand
\begin{equation}
  f_R = \frac{\Exp\!\left[w\,\Ind\{x \in R\}\right]}{\Exp[w]},
\end{equation}
estimated by
\begin{equation}
  \hat{f}_R(\mathcal{S})
  = \frac{\sum_{i:\,x_i \in R} w_i}{\sum_{i} w_i}.
  \label{eq:fhat}
\end{equation}
\end{definition}

Equation~\eqref{eq:fhat} is a ratio of weighted sums within one sample and
therefore satisfies Eq.~\eqref{eq:invariance}. Taking $R$ to select a charge
gives the charge ratio
\begin{equation}
  \hat{\rho} = \frac{\sum_{i:\, q_i = +1} w_i}{\sum_{i:\, q_i = -1} w_i},
\end{equation}
likewise scale invariant. The default region set consists of one-sided cuts on
the observables of Eq.~\eqref{eq:obs}, optionally restricted to one charge.

\subsection{Uncertainties: nonparametric bootstrap over units}
\label{sec:bootstrap}

Let $\mathcal{S}$ have units $1,\dots,m$ and let $\hat{\phi}(\mathcal{S})$ be
any of the above estimators. A bootstrap replica is formed by drawing $m$ unit
indices with replacement from $\{1,\dots,m\}$ and recomputing $\hat{\phi}$ from
the union of the drawn units~\cite{Efron1979,DavisonHinkley1997}. With $B$
replicas $\hat{\phi}^{*(1)},\dots,\hat{\phi}^{*(B)}$ we take
\begin{equation}
  \hat{\sigma}^{2}
  = \frac{1}{B-1}\sum_{b=1}^{B}
    \left(\hat{\phi}^{*(b)} - \overline{\hat{\phi}^{*}}\right)^{2},
  \qquad B = 400 .
\end{equation}

\begin{assumption}[Unit independence]
\label{ass:units}
The units of a sample are independent and identically distributed draws from a
population of units.
\end{assumption}

Assumption~\ref{ass:units} is what makes the bootstrap consistent here. It
justifies the choice of the file as the unit whenever enough files exist: files
correspond to independent simulation jobs, so their mutual independence is
plausible by construction, whereas events within a job share a configuration.
Resampling files propagates job-to-job variation into $\hat{\sigma}$
automatically, and in particular removes the need for the ad-hoc
error-inflation factors used in earlier analyses of these samples.

\subsection{Comparison statistic}

For a quantity $\phi$ measured in both samples we report the relative
difference and the standardised difference (``pull''),
\begin{equation}
  \Delta = \frac{\hat{\phi}_\Tt - \hat{\phi}_\Rr}{\hat{\phi}_\Rr},
  \qquad
  t = \frac{\hat{\phi}_\Tt - \hat{\phi}_\Rr}
           {\sqrt{\hat{\sigma}_\Tt^{2} + \hat{\sigma}_\Rr^{2}}} .
  \label{eq:pull}
\end{equation}
Addition of the variances in quadrature is valid because $\Rr$ and $\Tt$ are
independent samples.

\begin{remark}[Status of the pull]
\label{rem:pull}
Equation~\eqref{eq:pull} is a Wald-type statistic. Interpreting $t$ as a
standard normal deviate requires the sampling distribution of
$\hat{\phi}_\Tt - \hat{\phi}_\Rr$ to be approximately Gaussian, which holds
asymptotically in the number of units but is not guaranteed at $m_\Tt = 62$,
particularly if the weight is concentrated in a few units. Where a calibrated
interval is required, a studentised or bias-corrected accelerated bootstrap
interval~\cite{DavisonHinkley1997} should be substituted. We therefore quote
$\Delta$ with $\hat{\sigma}$ as the primary result and treat $t$ as an ordering
device.
\end{remark}

A quantity is reported as non-comparable, and excluded from all summaries, when
a region is empty in either sample or when $\hat{\sigma}/|\hat{\phi}| > 1/2$ in
either sample. Such a quantity is an absence of measurement, not a null result.

% ===========================================================================
\section{The two-sample Anderson--Darling test}
\label{sec:ad}

The tests of the previous section respond to a difference only where a cut is
placed. This section defines a test of the full one-dimensional distribution.

\subsection{Weighted empirical distribution functions}

\begin{definition}[Weighted EDF]
For a sample $\mathcal{S}$ and a scalar observable $x^{(j)}$,
\begin{equation}
  \hat{F}(t) = \frac{\sum_{i} w_i\, \Ind\{x^{(j)}_i \le t\}}{\sum_{i} w_i},
  \label{eq:edf}
\end{equation}
where $\Ind\{\cdot\}$ denotes the indicator function. $\hat{F}$ is
non-decreasing, right-continuous, and satisfies $\hat{F}(-\infty)=0$ and
$\hat{F}(+\infty)=1$.
\end{definition}

The normalisation in Eq.~\eqref{eq:edf} is precisely Eq.~\eqref{eq:invariance}:
$\hat{F}$ is unchanged under $w \mapsto aw$. The quantity $\hat{F}(t)$ is the
fraction of the sample's flux lying at or below $t$.

\subsection{The statistic}

Write $\hat{F}_\Rr, \hat{F}_\Tt$ for the two weighted EDFs and let
$\nu_\Rr, \nu_\Tt > 0$ be weights assigned to the samples (specified in
Sec.~\ref{sec:neffsub}), with $N = \nu_\Rr + \nu_\Tt$. Define the pooled
distribution function
\begin{equation}
  \hat{H}(t) = \frac{\nu_\Rr \hat{F}_\Rr(t) + \nu_\Tt \hat{F}_\Tt(t)}{N}.
\end{equation}

\begin{definition}[Two-sample Anderson--Darling statistic]
\begin{equation}
  A^{2}
  = \frac{\nu_\Rr \nu_\Tt}{N}
    \int_{-\infty}^{\infty}
    \frac{\bigl[\hat{F}_\Rr(t)-\hat{F}_\Tt(t)\bigr]^{2}}
         {\hat{H}(t)\bigl[1-\hat{H}(t)\bigr]}\ \dif \hat{H}(t).
  \label{eq:ad}
\end{equation}
\end{definition}

Equation~\eqref{eq:ad} is the two-sample form of the Anderson--Darling
statistic~\cite{AndersonDarling1952,AndersonDarling1954}, in the version due to
Pettitt~\cite{Pettitt1976} and generalised to $k$ samples by Scholz and
Stephens~\cite{ScholzStephens1987}. It belongs to the Cram\'er--von~Mises
family~\cite{Cramer1928,vonMises1928}
\begin{equation}
  \int \bigl[\hat{F}_\Rr - \hat{F}_\Tt\bigr]^{2}\, \psi\!\left(\hat{H}\right)
  \dif\hat{H},
  \label{eq:family}
\end{equation}
with weight function $\psi(H) = \{H(1-H)\}^{-1}$ in Eq.~\eqref{eq:family}; the choice $\psi \equiv 1$
gives the Cram\'er--von~Mises statistic $\omega^{2}$, and the supremum
$\sup_t|\hat{F}_\Rr - \hat{F}_\Tt|$ gives the Kolmogorov--Smirnov
statistic~\cite{Kolmogorov1933,Smirnov1939}. All three are computed and
reported.

\subsection{Interpretation of the weight function}
\label{sec:weighting}

The choice $\psi(H) = \{H(1-H)\}^{-1}$ is not an arbitrary emphasis of the
tails. For an EDF built from $N$ independent observations, $N\hat{H}(t)$ is
binomially distributed with mean $NH(t)$, so
\begin{equation}
  \Var\bigl[\hat{H}(t)\bigr] = \frac{H(t)\left[1-H(t)\right]}{N}.
\end{equation}
Equation~\eqref{eq:ad} therefore divides each squared discrepancy by its own
expected variance, making $A^2$ an integral of squared standardised residuals.
In the bulk $H \approx 1/2$ and $H(1-H) \approx 1/4$; at $H = 10^{-3}$ the
denominator is smaller by a factor $\approx 250$, so a discrepancy of given
absolute size in that region contributes $\approx 250$ times more. This is the
behaviour required by (P3): in a sparsely populated tail the two EDFs are
constrained to lie close together, so a separation there is correspondingly more
informative.

Reporting $A^2$, $\omega^2$ and $D$ together is itself diagnostic: rejection by
$A^2$ alone localises the discrepancy to the tails, whereas concordant rejection
by all three indicates a difference in the bulk.

\subsection{Distribution-freeness}

\begin{proposition}
\label{prop:df}
If $\hat{F}_\Rr$ and $\hat{F}_\Tt$ are the EDFs of independent samples of
i.i.d.\ observations from a common continuous distribution $F$, the law of
$A^{2}$ in Eq.~\eqref{eq:ad} does not depend on $F$.
\end{proposition}

\begin{proof}[Proof sketch]
Substituting $t \mapsto F^{-1}(s)$ maps the integral to one over $s \in (0,1)$
in which $F$ enters only through the probability integral transform
$F(x_i) \sim \mathrm{Uniform}(0,1)$. See Ref.~\cite{Pettitt1976}.
\end{proof}

\begin{remark}
Proposition~\ref{prop:df} concerns the \emph{two-sample} statistic. It is
distinct from the \emph{one-sample} Anderson--Darling goodness-of-fit test, for
which critical values depend on the hypothesised family and on whether its
parameters were estimated from the data~\cite{Stephens1974}. The frequent
objection that ``the Anderson--Darling test assumes normality'' refers to the
one-sample case and does not apply here.
\end{remark}

In this work Proposition~\ref{prop:df} is used only for orientation: no
tabulated critical values are employed, and the null distribution is constructed
by permutation (Sec.~\ref{sec:perm}), which requires neither continuity, nor
i.i.d.\ sampling, nor unit weights.

\subsection{Substitution of effective sample sizes}
\label{sec:neffsub}

For unweighted samples Eq.~\eqref{eq:ad} is defined with $\nu_\Rr = n_\Rr$ and
$\nu_\Tt = n_\Tt$. With weighted, correlated observations neither the counts nor
the binomial variance argument of Sec.~\ref{sec:weighting} applies exactly. We
set $\nu_\Rr = \Neff(\Rr)$ and $\nu_\Tt = \Neff(\Tt)$ from
Eq.~\eqref{eq:neff}.

This substitution is a heuristic: it is exact only for equally weighted
independent observations, and Proposition~\ref{prop:split} shows it to be biased
for split samples. Its consequences are nonetheless benign, for the following
reason.

\begin{proposition}[Validity is independent of the choice of statistic]
\label{prop:valid}
Let $S$ be \emph{any} real-valued function of the labelled data, and let its
null distribution be obtained by permutation of unit labels as in
Sec.~\ref{sec:perm}. Then under $H_0$ of Definition~\ref{def:h0} the resulting
p-value is valid, i.e.\ $\Prob(\hat{p} \le \alpha) \le \alpha$ for all
$\alpha \in (0,1)$.
\end{proposition}

\begin{proof}[Proof sketch]
Under exchangeability the observed labelling is, conditionally on the multiset
of units, uniformly distributed over the group of size-preserving relabellings.
The observed value of $S$ is therefore exchangeable with the permuted values,
and the rank-based p-value of Eq.~\eqref{eq:pval} is super-uniform. See
Refs.~\cite{Fisher1935,Pitman1937,LehmannRomano2005}.
\end{proof}

Proposition~\ref{prop:valid} is the central structural point. The prefactor
$\nu_\Rr\nu_\Tt/N$ in Eq.~\eqref{eq:ad} is a fixed positive constant within a
given permutation experiment, and the weighting $\{H(1-H)\}^{-1}$ is a fixed
function of the pooled data; neither affects the exchangeability argument.
Approximations in $\nu$ therefore change the \emph{power} of the test --- which
alternatives it detects most readily --- but not the validity of its p-values.
The same argument covers the discretisation of Sec.~\ref{sec:discrete}.

\subsection{The discrete form evaluated by the implementation}
\label{sec:discrete}

The implementation does not evaluate Eq.~\eqref{eq:ad} on the exact EDFs. Both
samples are binned onto a common grid of $K$ contiguous bins with edges
$t_0 < t_1 < \dots < t_K$, chosen so that each bin carries an equal share of the
pooled weight (Sec.~\ref{sec:binning}). Writing $h^\Rr_k, h^\Tt_k$ for the
weighted contents of bin $k$ and
\begin{equation}
  \hat{F}^{\Rr}_k = \frac{\sum_{l \le k} h^\Rr_l}{\sum_{l} h^\Rr_l},
  \qquad
  \hat{F}^{\Tt}_k = \frac{\sum_{l \le k} h^\Tt_l}{\sum_{l} h^\Tt_l},
  \qquad
  \hat{H}_k = \frac{\nu_\Rr \hat{F}^{\Rr}_k + \nu_\Tt \hat{F}^{\Tt}_k}{N},
\end{equation}
with $\hat{H}_0 = 0$ and $\Delta\hat{H}_k = \hat{H}_k - \hat{H}_{k-1}$, the
computed statistics are
\begin{align}
  A^{2} &= \frac{\nu_\Rr\nu_\Tt}{N} \sum_{k=1}^{K-1}
           \frac{\bigl(\hat{F}^{\Rr}_k-\hat{F}^{\Tt}_k\bigr)^{2}}
                {\hat{H}_k\bigl(1-\hat{H}_k\bigr)}\,\Delta \hat{H}_k ,
           \label{eq:addisc}\\
  \omega^{2} &= \frac{\nu_\Rr\nu_\Tt}{N} \sum_{k=1}^{K}
           \bigl(\hat{F}^{\Rr}_k-\hat{F}^{\Tt}_k\bigr)^{2}\,\Delta \hat{H}_k ,
           \qquad
  D = \sqrt{\frac{\nu_\Rr\nu_\Tt}{N}}\,
      \max_{k}\bigl|\hat{F}^{\Rr}_k-\hat{F}^{\Tt}_k\bigr| .
\end{align}
The term $k=K$ is excluded from Eq.~\eqref{eq:addisc} because $\hat{H}_K = 1$
exactly, at which the integrand is singular; this is the discrete counterpart of
the convergence of the integral in Eq.~\eqref{eq:ad}. An optional threshold
$\epsilon$ may be imposed, discarding terms with
$\hat{H}_k(1-\hat{H}_k) \le \epsilon$, to limit the influence of extreme-weight
observations near the boundaries.

By Proposition~\ref{prop:valid}, Eq.~\eqref{eq:addisc} defines a valid test
statistic irrespective of $K$. The dependence of the \emph{value} of $A^2$ on
$K$ is measured directly and reported (Sec.~\ref{sec:binstab}).

\subsection{Binning}
\label{sec:binning}

Bin edges are placed at equal quantiles of the pooled weighted distribution, so
that each of the $K$ bins carries $1/K$ of the pooled weight. Two properties
motivate this choice. First, it equalises the statistical precision across bins,
so that sparsely populated tails are represented by wide bins containing usable
statistics rather than by many near-empty narrow bins. Second, it gives
$\Delta\hat{H}_k \approx 1/K$, so that the sum in Eq.~\eqref{eq:addisc}
approximates the integral uniformly.

For display, densities $h_k / (t_k - t_{k-1})$ are plotted, since the bin
contents are flat by construction.

% ===========================================================================
\section{Calibration by permutation}
\label{sec:perm}

Let $S$ denote any of $A^2$, $\omega^2$, $D$. Let $\pi$ range over relabellings
of the $m_\Rr + m_\Tt$ units that preserve the group sizes, and let $S^{(\pi)}$
be the statistic recomputed under $\pi$, with each pseudo-sample renormalised
exactly as the observed samples are. Drawing $M$ relabellings uniformly at
random, the p-value is
\begin{equation}
  \hat{p} = \frac{1 + \#\left\{\pi : S^{(\pi)} \ge S^{\mathrm{obs}}\right\}}{1+M}.
  \label{eq:pval}
\end{equation}

The test is one-sided because all three statistics are non-negative and large
values indicate discrepancy. The addition of $1$ to numerator and denominator is
required for validity when $M$ is finite~\cite{PhipsonSmyth2010}; it implies the
attainable minimum $\hat{p} = 1/(1+M)$, and results at this bound are flagged as
such rather than reported as smaller.

Two properties are essential. Units are permuted as wholes, so that within-unit
dependence --- including splitting copies and muons of a common primary --- is
preserved under $H_0$ and does not narrow the null distribution. And each
replica is renormalised, so the null is a pure shape null: no relabelling can
respond to the normalisation difference of Assumption~\ref{ass:scale}.

\begin{remark}[The reported $z$ is not a significance]
\label{rem:z}
The implementation also reports
\begin{equation*}
  z = \frac{S^{\mathrm{obs}} - \overline{S^{(\pi)}}}
           {\mathrm{sd}\bigl(S^{(\pi)}\bigr)} .
\end{equation*}
The permutation null of $A^2$ is strongly
right-skewed, so $z$ admits no interpretation as a normal deviate. It is
retained solely to order variables whose p-values all attain the bound of
Eq.~\eqref{eq:pval}; the tables of Sec.~\ref{sec:results} are read accordingly.
\end{remark}

\subsection{Localisation of the discrepancy}
\label{sec:profile}

A p-value states whether the samples differ, not where. Since $A^2$ is an
integral, its integrand is reported. For each observable we plot

\begin{enumerate}
\item the signed difference $\hat{F}_\Tt(t) - \hat{F}_\Rr(t)$, with pointwise
$68\%$ and $95\%$ envelopes formed from the empirical quantiles of
$\hat{F}^{(\pi)}_\Tt - \hat{F}^{(\pi)}_\Rr$ over the same $M$ relabellings.
Positive values indicate that $\Tt$ has accumulated a larger share of its flux
below $t$;
\item the normalised partial sums of Eq.~\eqref{eq:addisc},
$C_k = \sum_{l \le k} a_l \big/ \sum_{l} a_l$, where $a_l$ denotes the $l$-th
summand, whose increments identify the bins generating the statistic.
\end{enumerate}

The second localises better than the first. Because $\hat{F}_\Tt - \hat{F}_\Rr$
is a difference of cumulative functions constrained to vanish at both ends, a
discrepancy concentrated at one value appears as a broad excursion; $C_k$ is a
density in $k$ and exhibits a step. A steady rise of $C_k$ indicates a shift
distributed over the range, a step indicates a localised feature.

\begin{remark}[Pointwise, not simultaneous]
The envelopes in item~1 are pointwise quantiles. Under $H_0$ approximately
$5\%$ of the $K$ bins lie outside the $95\%$ envelope, so inference must be
drawn from coherent runs of bins. No simultaneous confidence band is
constructed.
\end{remark}

% ===========================================================================
\section{Implementation}

\subsection{Streaming accumulation}

A single production file contains up to $3.6\times10^{7}$ selected muons and a
full sample up to $2.3\times10^{9}$; per-muon storage is therefore infeasible.
The implementation accumulates muons directly into fixed histograms indexed by
(unit, charge, observable) with $n_{\mathrm{fine}}$ bins per observable, so that
memory scales as
\begin{equation}
  \mathcal{O}\!\left(m \times 2 \times d \times n_{\mathrm{fine}}\right),
\end{equation}
independent of the number of muons. All files of a sample are processed as a
single dataset with implicit multi-threading using
\texttt{ROOT::RDataFrame}~\cite{ROOT,RDataFrame}. All statistics of
Secs.~\ref{sec:fractions}--\ref{sec:perm} are subsequently computed from these
histograms; the test bins of Sec.~\ref{sec:binning} are formed by merging fine
bins.

\begin{proposition}[Exactness of merging]
\label{prop:merge}
If the coarse bin edges form a subset of the fine bin edges, the coarse bin
contents obtained by summing fine bins are identical to those obtained by
accumulating directly on the coarse grid.
\end{proposition}

Proposition~\ref{prop:merge} is what permits the binning-stability scan of
Sec.~\ref{sec:binstab} to be performed without re-reading the data. It is
verified numerically to $1.8\times10^{-15}$.

\subsection{Approximations introduced}

Two approximations follow from accumulating on a fixed grid, both quantified at
runtime.

\begin{enumerate}
\item \emph{Cut snapping.} A region boundary is displaced to the nearest fine
bin edge, by at most half a fine bin width. With $n_{\mathrm{fine}} = 2048$ over
the momentum range this is $\lesssim 0.03\%$ in $p$.
\item \emph{Clamping.} Values outside the recorded range are accumulated in the
end bins. The recorded range is fixed by a probe over the leading entries of the
first files of each sample, padded by $35\%$. The clamped fraction is reported;
it was below $2\times10^{-6}$ in all runs of Sec.~\ref{sec:results}.
\end{enumerate}

Neither affects the validity of the permutation p-values, by
Proposition~\ref{prop:valid}.

% ===========================================================================
\section{Validation}
\label{sec:validation}

\subsection{Unit tests of the statistical core}

The numerical core is free of external dependencies and is covered by nine
assertions on synthetic data, of which the substantive ones are:

\begin{center}
\small
\begin{tabular}{@{}rp{0.78\textwidth}@{}}
\toprule
\# & Assertion \\
\midrule
1 & Multiplying one sample's weights by $10^{3}$ leaves $A^{2}$, $\omega^{2}$
    and $\hat{p}$ bit-identical, verifying Eq.~\eqref{eq:invariance}. \\
2 & Two independent samples from a common law yield p-values that are not
    systematically small. \\
3 & For two laws with identical mean \emph{and} identical variance but
    different tail behaviour, a Fisher-style standardised mean separation
    returns $+0.0002$ while $A^{2}$ attains the bound of Eq.~\eqref{eq:pval}
    with $z = +176$. This demonstrates the failure mode asserted in (P2). \\
4 & An injected change of the charge ratio from $1.00$ to $1.50$ is detected at
    $t = 11$ while an unchanged fraction is not, in the presence of a
    $500\times$ normalisation difference. \\
5 & Unit-level permutation is not anti-conservative in the presence of exact
    duplicate observations within a unit. \\
7 & The streaming implementation reproduces the exact per-muon evaluation of
    $A^{2}$ to $0.21\%$ with identical p-values. \\
9 & Merging of fine bins is exact to $1.8\times10^{-15}$
    (Proposition~\ref{prop:merge}). \\
\bottomrule
\end{tabular}
\end{center}

\subsection{Closure}

Supplying the same file as both inputs yields $A^{2} \sim 10^{-27}$ (floating
point zero), $\hat{p} = 1$ and $t = 0$ for every quantity. Partitioning a single
sample into two halves by unit and comparing them yields no pull exceeding $3$
and p-values distributed over $(0,1)$.

\subsection{Calibration of the p-value distribution}

Closure demonstrates the absence of spurious differences but not the calibration
of $\hat{p}$. Under $H_0$, $\hat{p}$ must be super-uniform on $(0,1)$. The
half-splitting procedure was therefore repeated $20$ times with independent
random partitions, yielding $20$ independent p-values per observable. The
observed means ranged over $[0.427, 0.546]$ against an expected $0.500$ with
standard error $0.065$; the fraction below $0.05$ was consistent with $0.05$;
and the Kolmogorov distance from uniformity did not exceed the $5\%$ critical
value for $n=20$. A P--P plot of the ordered p-values against $k/(n+1)$ is
produced.

This procedure resolved an apparent anomaly: across three individual runs,
seven or eight of eight $z$ values were negative, suggesting conservatism. The
scan showed the p-values to be calibrated. The apparent pattern arises because
the observables of Eq.~\eqref{eq:obs} are strongly dependent --- four are
functions of the same momentum components --- so a single run furnishes
approximately two independent p-values rather than eight.

\subsection{Binning stability}
\label{sec:binstab}

By Proposition~\ref{prop:merge}, $A^{2}$ can be recomputed on grids coarser by
factors of $4$ and $16$ from a single pass. Any observable whose value changes
by more than $35\%$ across these grids is flagged as unconverged.

This check was introduced in response to an observed failure. With
$n_{\mathrm{fine}} = 512$, the observable $x^{(5)} = x$ returned $A^{2} = 13.3$
at $\hat{p} < 0.002$, having been compatible in all previous comparisons. Its
profile (Sec.~\ref{sec:profile}) showed $C_k$ rising as a step over
$\approx 10$~cm at $x \approx 0$, the signature of a sharp feature straddling a
bin edge. Refinement to $n_{\mathrm{fine}} = 2048$ reduced $A^{2}$ to $4.5$ and
to $2.29$ at $8192$, where the three grids agree to better than $1\%$. A genuine
discrepancy does not diminish under refinement. The default was accordingly
raised to $2048$.

\subsection{Audit of the weight bookkeeping of split productions}
\label{sec:audit}

A production employing splitting writes each muon $k$ times with weight $w/k$.
If performed correctly this is a variance-reduction transformation leaving all
weighted distributions invariant. If the division is omitted or incorrect, every
weighted comparison against an unsplit sample is biased, and no property of the
statistical procedure can compensate.

Splitting acts within a single tree entry, which yields an invariant requiring
no duplicate detection:
\begin{equation}
  \sum_{i \in e} w_i \ \text{is unchanged by splitting for every entry } e
  \quad\Longrightarrow\quad
  \frac{\Exp_\Tt\!\left[\sum_{i \in e} w_i\right]}
       {\Exp_\Rr\!\left[\sum_{i \in e} w_i\right]} = 1 .
  \label{eq:auditinv}
\end{equation}
Two independent estimates of $k$ are also available: from the ratio of mean
multiplicities, and from the inverse ratio of mean per-muon weights. For the
samples of Sec.~\ref{sec:results} these gave $k = 64.5$ and $k = 69.2$,
agreeing to $7.3\%$, with the ratio in Eq.~\eqref{eq:auditinv} equal to
$0.932$. Absence of the division would have produced $\approx 65$ in place of
$\approx 1$.

\begin{remark}
Equation~\eqref{eq:auditinv} constrains the bookkeeping in aggregate. A
misallocation among the copies of a single muon that preserves the per-entry sum
would satisfy it. Establishing per-muon correctness requires duplicate
identification and was not performed.
\end{remark}

\begin{remark}[A sample-selection hazard]
Comparisons using a different weight branch of the same jobs produced an
apparent $50\%$ difference in the momentum spectrum: the same muons gave
$\sum_i w_i = 5.36\times10^{5}$ against $1.78\times10^{4}$. Consistency of
branch selection between samples is a precondition for all of the above.
\end{remark}

% ===========================================================================
\section{Application}
\label{sec:results}

\subsection{Staged comparison}

The procedure was applied in stages of increasing scope, each providing a
control for the next. All entries are $A^{2}$ for $x^{(1)} = \log_{10} p$.

\begin{center}
\begin{tabular}{@{}llrl@{}}
\toprule
Stage & Comparison & $A^{2}$ & $\hat{p}$ \\
\midrule
1 & one file against itself & $\sim 10^{-27}$ & $1$ \\
2 & two jobs of the reference production & $0.25$ & $0.96$ \\
3 & one file from each production & $7.6$ & $<0.002$ \\
4 & complete file lists ($m_\Rr = 205$, $m_\Tt = 62$) & $569$ & $<0.0005$ \\
\bottomrule
\end{tabular}
\end{center}

Stage~2 is the essential control: two jobs of a single production, differing
only in random seed and proton count, are compatible in the observable that
separates the productions by more than three orders of magnitude in $A^{2}$ at
stage~4.

\subsection{Complete comparison}

At stage~4, with $2.33\times10^{9}$ test-sample muons over $62$ units and
$m_\Rr = 205$ reference units, every observable and every charge selection
rejects at the bound $\hat{p} < 5\times10^{-4}$. At this sample size the
p-value has ceased to discriminate, since differences far below any level of
physical relevance are resolvable. We therefore report effect sizes and the
ordering of $A^{2}$.

\begin{center}
\begin{tabular}{@{}lrrr@{}}
\toprule
Observable & $A^{2}$ (both) & $A^{2}$ ($\mu^+$) & $A^{2}$ ($\mu^-$) \\
\midrule
$\log_{10} p$              & $569.3$ & $226.8$ & $79.0$ \\
$\log_{10} p_z$            & $569.7$ & $227.1$ & $79.0$ \\
$r$                        & $324.2$ & $117.3$ & $52.6$ \\
$\log_{10}\theta$          & $245.3$ & $96.7$  & $34.9$ \\
$y$                        & $64.7$  & $24.0$  & $10.4$ \\
$x$                        & $64.2$  & $22.8$  & $10.8$ \\
$q\,x$                     & $63.4$  & $22.8$  & $10.7$ \\
$\operatorname{arcsinh} p_T$ & $23.3$ & $10.1$  & $3.0$ \\
\bottomrule
\end{tabular}
\end{center}

\begin{center}
\begin{tabular}{@{}lrrr@{}}
\toprule
Quantity & $\hat{\phi}_\Rr$ & $\hat{\phi}_\Tt$ & $\Delta$ \\
\midrule
$f(p>50~\mathrm{GeV})$        & $0.04861 \pm 0.00006$ & $0.04961 \pm 0.00012$ & $+2.06\%$ \\
$f(p>100~\mathrm{GeV})$       & $0.00842 \pm 0.00003$ & $0.00858 \pm 0.00004$ & $+1.91\%$ \\
$f(p_T>1~\mathrm{GeV})$       & $0.02674 \pm 0.00006$ & $0.02717 \pm 0.00008$ & $+1.60\%$ \\
$f(\theta>20~\mathrm{mrad})$  & $0.67995 \pm 0.00019$ & $0.67649 \pm 0.00020$ & $-0.51\%$ \\
$f(\theta>50~\mathrm{mrad})$  & $0.23674 \pm 0.00016$ & $0.23346 \pm 0.00017$ & $-1.39\%$ \\
$\rho = N(\mu^+)/N(\mu^-)$    & $1.13765 \pm 0.00092$ & $1.12521 \pm 0.00130$ & $-1.09\%$ \\
\bottomrule
\end{tabular}
\end{center}

Uncertainties are bootstrap standard deviations over file units; see
Remark~\ref{rem:pull} concerning their interpretation as normal deviates.

\subsection{Interpretation}

The ordering of $A^{2}$ does not constitute a set of independent findings. Since
$\theta = \arctan(p_T/|p_z|)$ and $p_T$ ranks lowest, by a factor $24$ relative
to $p$, the discrepancy in $\theta$ is consistent with being induced by that in
$p$; and $r$, the radial coordinate at a fixed longitudinal position, follows
from $\theta$. The observables $x$, $y$ and $q\,x$ rank an order of magnitude
below $p$.

The profile of $\log_{10} p$ shows a single-signed excursion of
$\hat{F}_\Tt - \hat{F}_\Rr$ extending over the bulk of the spectrum, outside the
$95\%$ envelope across several hundred consecutive bins, with $C_k$ rising
monotonically and without a step. The discrepancy is therefore distributed over
the spectrum rather than localised, and the absence of a sign change indicates
that the positive excursions observed in the highest-momentum fractions are
statistical fluctuations superposed on a monotone trend.

In summary: the test production exhibits a harder muon momentum spectrum than
the reference, by approximately $2\%$ in the fraction above $50$~GeV, with
associated narrowing of the angular and radial distributions; the charge ratio
is lower by approximately $1\%$; and all remaining observables differ by less.
Whether differences of this magnitude are material is a question about the
downstream application, not about the present analysis.

\subsection{Limitations}

\begin{enumerate}
\item \textbf{Marginals only.} The procedure tests $d$ one-dimensional marginal
distributions. An alternative in which the marginals agree but the dependence
structure differs is not detected. A multivariate extension --- for instance an
energy or kernel two-sample statistic~\cite{Szekely2004,Gretton2012}, or a
classifier-based test~\cite{Friedman2003,Lopez2017} --- would address this and is
not implemented.

\item \textbf{Multiplicity.} Twenty-four tests ($8$ observables $\times$ $3$
charge selections) are reported without correction. The observables are strongly
dependent, so neither a Bonferroni bound nor an assumption of independence is
appropriate; no valid joint statement about the family is made.

\item \textbf{The pull is not a calibrated significance.} See
Remark~\ref{rem:pull}. In particular, the number of effectively independent
units may be smaller than $m$ if the weight is concentrated in few units, in
which case $\hat{\sigma}$ is underestimated. This has not been quantified: the
per-unit weight distribution should be examined and the effective number of
units, $\left(\sum_f S_f\right)^2 \big/ \sum_f S_f^2$ with $S_f$ the total
weight of unit $f$, reported.

\item \textbf{Reproducibility of intermediate runs.} Two subsampled runs over
the same file lists produced $\Delta = +2.76\%$ and $+0.07\%$ for
$f(p>50~\mathrm{GeV})$, a spread exceeding the quoted uncertainties. As the
charge ratio, which involves no cut, was likewise affected, this cannot be
attributed to binning. The cause was not established. The complete-statistics
results of this section are internally consistent, but the discrepancy is
recorded as unresolved.

\item \textbf{The geometric interpretation is an inference.} The claim that the
$\theta$ and $r$ discrepancies are induced by that in $p$ rests on the ordering
of $A^{2}$ and on the comparatively small discrepancy in $p_T$; it is not a
quantitative demonstration. Reweighting $\Rr$ by the estimated ratio of momentum
densities and retesting $\theta$ would provide one.

\item \textbf{$\Neff$ for split samples.} By Proposition~\ref{prop:split},
reported values of $\Neff$ for the test sample are upper bounds and must not be
used to infer precision.
\end{enumerate}

% ===========================================================================
\section{Summary}

Comparison of weighted samples whose normalisation is arbitrary requires
statistics invariant under independent rescaling of either sample, which
excludes any procedure operating on absolute rates. Within that constraint, the
weighted two-sample Anderson--Darling statistic calibrated by permutation of
resampling units provides a test that is valid under the exchangeability
hypothesis of Definition~\ref{def:h0} without assumptions on the underlying
distributions, and whose weight function assigns to each region of the
observable a contribution inversely proportional to its variance, as required
when tail behaviour is of interest. Ratios of weighted sums within a sample,
with uncertainties from a bootstrap over units, supply the complementary
comparison of relative amounts with correlation-aware uncertainties.

Proposition~\ref{prop:valid} establishes that the two principal approximations
of the implementation --- the substitution of effective sample sizes for counts,
and the discretisation of the empirical distribution functions --- affect the
power of the procedure but not the validity of its p-values. The remaining
approximations are quantified at runtime.

\begin{thebibliography}{99}

\bibitem{AndersonDarling1952}
T.~W. Anderson and D.~A. Darling,
\emph{Asymptotic theory of certain ``goodness of fit'' criteria based on
stochastic processes},
Ann.\ Math.\ Statist.\ \textbf{23} (1952) 193--212.

\bibitem{AndersonDarling1954}
T.~W. Anderson and D.~A. Darling,
\emph{A test of goodness of fit},
J.\ Amer.\ Statist.\ Assoc.\ \textbf{49} (1954) 765--769.

\bibitem{Pettitt1976}
A.~N. Pettitt,
\emph{A two-sample Anderson--Darling rank statistic},
Biometrika \textbf{63} (1976) 161--168.

\bibitem{ScholzStephens1987}
F.~W. Scholz and M.~A. Stephens,
\emph{$k$-sample Anderson--Darling tests},
J.\ Amer.\ Statist.\ Assoc.\ \textbf{82} (1987) 918--924.

\bibitem{Stephens1974}
M.~A. Stephens,
\emph{EDF statistics for goodness of fit and some comparisons},
J.\ Amer.\ Statist.\ Assoc.\ \textbf{69} (1974) 730--737.

\bibitem{Cramer1928}
H.~Cram\'er,
\emph{On the composition of elementary errors},
Skand.\ Aktuarietidskr.\ \textbf{11} (1928) 141--180.

\bibitem{vonMises1928}
R.~von Mises,
\emph{Wahrscheinlichkeit, Statistik und Wahrheit}, Springer, 1928.

\bibitem{Kolmogorov1933}
A.~N. Kolmogorov,
\emph{Sulla determinazione empirica di una legge di distribuzione},
Giorn.\ Ist.\ Ital.\ Attuari \textbf{4} (1933) 83--91.

\bibitem{Smirnov1939}
N.~V. Smirnov,
\emph{On the estimation of the discrepancy between empirical curves of
distribution for two independent samples},
Bull.\ Math.\ Univ.\ Moscou \textbf{2} (1939) 3--14.

\bibitem{Fisher1935}
R.~A. Fisher, \emph{The Design of Experiments}, Oliver and Boyd, 1935.

\bibitem{Pitman1937}
E.~J.~G. Pitman,
\emph{Significance tests which may be applied to samples from any populations},
Suppl.\ J.\ Roy.\ Statist.\ Soc.\ \textbf{4} (1937) 119--130.

\bibitem{LehmannRomano2005}
E.~L. Lehmann and J.~P. Romano,
\emph{Testing Statistical Hypotheses}, 3rd ed., Springer, 2005, Ch.~15.

\bibitem{PhipsonSmyth2010}
B.~Phipson and G.~K. Smyth,
\emph{Permutation p-values should never be zero: calculating exact p-values when
permutations are randomly drawn},
Stat.\ Appl.\ Genet.\ Mol.\ Biol.\ \textbf{9} (2010) Article~39.

\bibitem{Kish1965}
L.~Kish, \emph{Survey Sampling}, Wiley, 1965, Sec.~11.7.

\bibitem{Efron1979}
B.~Efron,
\emph{Bootstrap methods: another look at the jackknife},
Ann.\ Statist.\ \textbf{7} (1979) 1--26.

\bibitem{DavisonHinkley1997}
A.~C. Davison and D.~V. Hinkley,
\emph{Bootstrap Methods and their Application}, Cambridge University Press,
1997.

\bibitem{Szekely2004}
G.~J. Sz\'ekely and M.~L. Rizzo,
\emph{Testing for equal distributions in high dimension},
InterStat, November 2004.

\bibitem{Gretton2012}
A.~Gretton, K.~M. Borgwardt, M.~J. Rasch, B.~Sch\"olkopf and A.~Smola,
\emph{A kernel two-sample test},
J.\ Mach.\ Learn.\ Res.\ \textbf{13} (2012) 723--773.

\bibitem{Friedman2003}
J.~H. Friedman,
\emph{On multivariate goodness of fit and two sample testing},
eConf \textbf{C030908} (2003) THPD002.

\bibitem{Lopez2017}
D.~Lopez-Paz and M.~Oquab,
\emph{Revisiting classifier two-sample tests},
Proc.\ ICLR, 2017.

\bibitem{ROOT}
R.~Brun and F.~Rademakers,
\emph{ROOT --- an object oriented data analysis framework},
Nucl.\ Instrum.\ Meth.\ A \textbf{389} (1997) 81--86.

\bibitem{RDataFrame}
D.~Piparo, P.~Canal, E.~Guiraud \emph{et al.},
\emph{RDataFrame: easy parallel ROOT analysis at 100 threads},
EPJ Web Conf.\ \textbf{214} (2019) 06029.

\end{thebibliography}

\appendix
\section{Procedure}

\begin{lstlisting}[language=bash]
# build and run the unit tests of Sec. 8.1
cmake -S . -B build && cmake --build build -j && ctest --test-dir build

# closure (Sec. 8.2)
build/flux_compare REF.root --closure

# calibration of the p-value distribution (Sec. 8.3)
build/flux_compare REF.root --closure-scan 20 --nperm 200

# weight audit (Sec. 8.5)
build/flux_weight_audit REF.root TEST.root

# comparison; the first argument is the reference sample
build/flux_compare @ref.txt @test.txt \
    --label-ref nosplit --label-test split100 \
    --split-charge --nfine 2048 --nperm 2000 --outdir out
\end{lstlisting}

\noindent
For each observable and charge selection the procedure emits: the two densities
with a shape-ratio panel; the discrepancy profile of Sec.~\ref{sec:profile}; the
permutation null of Eq.~\eqref{eq:pval} with the observed value indicated; and
the per-bin effective sample size of Eq.~\eqref{eq:neff}.

\end{document}